\exercisesection{Extension of scalars in regular algebras}

\begin{enumerate}
\def\labelenumi{\arabic{enumi})}
\tightlist
\item
  Let \(k\) be a non-perfect field of characteristic \(p > 2\) , let \(a\) be an element of \(k - k^p\) , and let \(A = k[\mathrm{X},\mathrm{Y}] / (\mathrm{X}^2 -\mathrm{Y}^p +a)\) .
\end{enumerate}

\begin{enumerate}
\def\labelenumi{\alph{enumi})}
\item
  Prove that the ring A is regular (one may prove that A is integrally closed in its field of fractions \(K = k(Y)[X]/(X^{2} - Y^{p} + a)\) ).
\item
  Prove that k is algebraically closed in K.
\item
  Let \(k'\) the extension \(k(a^{1/p})\) of \(k\) prove that the ring \(A \otimes_k k'\) is not normal.
\end{enumerate}

\begin{enumerate}
\def\labelenumi{\arabic{enumi})}
\setcounter{enumi}{1}
\tightlist
\item
  Let \(k\) be a field, A an essentially finitely generated \(k\)-algebra, B a \(k\)-algebra, and \(p\) an integer.
\end{enumerate}

\begin{enumerate}
\def\labelenumi{\alph{enumi})}
\item
  If A and B satisfy the property \(S_{p}\) (§ 1, exerc. 8), then so does \(A \otimes_{k} B\) (argue as in the proof of prop. 4, using loc. cit., f)).
\item
  Let K be an extension of k; for \(\mathrm{A}_{(\mathbb{K})}\) check \(S_{p}\) it is necessary and sufficient that it be so for A.
\end{enumerate}

\begin{enumerate}
\def\labelenumi{\arabic{enumi})}
\setcounter{enumi}{2}
\tightlist
\item
  Let A be a Noetherian ring and \(\rho : A \to B\) a ring homomorphism. We say that the A-algebra B is absolutely regular if B is a Noetherian ring and a flat A-module and, for every \(\mathfrak{p} \in \operatorname{Spec}(A)\), the \(\kappa(\mathfrak{p}) \otimes_{A} B\) is absolutely regular.
\end{enumerate}

\begin{enumerate}
\def\labelenumi{\alph{enumi})}
\item
  Let T be a multiplicative subset of B and S a subset of A such that \(\rho(S) \subset T\) If B is a regular A-algebra, \(T^{1}B\) is an \(S^{-1}A\) regular algebra.
\item
  For a Noetherian A-algebra B to be absolutely regular, it is necessary and sufficient that for every maximal ideal n of B, the \(A_{\rho^{-1}(n)}\) -algebra \(B_{n}\) be absolutely regular.
\item
  If C is an absolutely regular B-algebra and B is an absolutely regular A-algebra, then C is an absolutely regular A-algebra.
\item
  If B is a Noetherian ring and C is a faithfully flat B-module and an absolutely regular A-algebra, then B is an absolutely regular A-algebra.
\item
  Let B be an absolutely regular A-algebra and A' an A-algebra essentially of finite type. The A'-algebra A' ⊗A B is absolutely regular.
\end{enumerate}

\begin{enumerate}
\def\labelenumi{\arabic{enumi})}
\setcounter{enumi}{3}
\item
  Let A be a Noetherian ring, B a faithfully flat absolutely regular A-algebra (exerc. 3). For A to be a regular (resp. normal, resp. Macaulay, resp. Gorenstein) ring, it is necessary and sufficient that the same be true of B.
\item
  One says that A is a Grothendieck ring if A is a Noetherian ring and, for every prime ideal p of A, the completion of the local ring \(\Lambda_{p}\) is an \(A_{p}\) -algebra absolutely regular (exerc. 3).
\end{enumerate}

\begin{enumerate}
\def\labelenumi{\alph{enumi})}
\item
  Prove that every ring of fractions and every quotient of a Grothendieck ring is a Grothendieck ring.
\item
  Let A be a Noetherian ring and B a faithfully flat absolutely regular A-algebra. If B is a Grothendieck ring, prove that the same is true of A (if p is a prime ideal of A and q a prime ideal of B above p, one will show that \(\widehat{B}_{q}\) is a \(\widehat{\Lambda_{p}}\) -module faithfully flat using III, § 5, n° 4, prop. 4; one will then apply exerc. 3).
\end{enumerate}

\begin{enumerate}
\def\labelenumi{\arabic{enumi})}
\setcounter{enumi}{5}
\item
  Let A be a Grothendieck ring, I an ideal of A, and \(\widehat{\mathsf{A}}\) the separated completion of A for the I-adic topology. Prove that \(\widehat{\mathsf{A}}\) is an absolutely regular A-algebra. (Let \(\mathfrak{n}\) be a maximal ideal of \(\widehat{\mathsf{A}}\) , and by \(\mathfrak{m}\) its inverse image in A; deduce from exerc. 3 that \(\widehat{\Lambda}_{\mathfrak{n}}\) is an \(\mathrm{A}_{\mathfrak{m}}\) -algebra absolutely regular.)
\item
  \begin{enumerate}
  \def\labelenumii{\alph{enumii})}
  \tightlist
  \item
    Show that an integral Dedekind ring of characteristic 0 is a Grothendieck ring.
  \end{enumerate}
\end{enumerate}

\begin{enumerate}
\def\labelenumi{\alph{enumi})}
\setcounter{enumi}{1}
\tightlist
\item
  Let \(k\) a field of characteristic \(p > 0\) such that \(k\) be of infinite dimension over \(k^p\) , \(r\) an integer \(\geqslant 1\) , and A the subring of \(k[[\mathrm{X}_1, \ldots, \mathrm{X}_r]]\) formed by formal series whose coefficients generate a vector space of finite dimension over \(k^p\) . Show that A is a regular local ring of dimension \(r\) whose completion is identified with \(k[[\mathrm{X}_1, \ldots, \mathrm{X}_r]]\) , but is not a Grothendieck ring (use exerc. 17 of III, § 3).
\end{enumerate}

